\section{Discussion}
\label{sec: discussion}

Our study reveals the diverse category learning trajectories humans exhibit in a multi-dimensional feature space, shaped by individualized information processing (Fig.~\ref{fig:3}). Beyond reproducing choices, we recover trial-wise internal strategies that evolve over time (Fig.~\ref{fig:4}), using a modular Bayesian framework (Fig.~\ref{fig:2}) that balances exploratory and exploitative feature strategies under cognitive constraints (Fig.~\ref{fig:5}). Our findings highlight the heterogeneity of human learning and the importance of modeling internal processes, not just behavior, to understand adaptive cognition.

The mathematical foundation of our model conceptualizes learning as a dynamic information optimizer, combining both global grid-based flexible sampling and local gradient-driven optimization. This dual mechanism manifests particular strength in perceptual decision-making, where our formulation permits a high degree of flexibility in the partitioning of the stimulus space. Such theoretical flexibility supports the construction of a hypothesis space encompassing infinite decision rules – ranging from linear boundaries to complex nonlinear configuration ~\cite{jaworska2022different} – as well as varying level of computational complexity (from simple 1-bit to high-dimensional multi-bit, ~\cite{cohen2013high}), and naturally extends to diverse conceptual structure (e.g., hierarchical, tree-like representations). By parameterizing task-specific internal states in a feature-based hypothesis space, our approach can bridge internal model dynamics with neural correlates of decision-making ~\cite{yang2007probabilistic}, probabilistic inference ~\cite{soltani2006neural}, and adaptive learning shifts \cite{nassar2010approximately}, thereby advancing cognitive algorithms reverse-engineering from brain activity.

\textbf{Limitation.} Verbal reports provided ground truth for participants’ internal representations but may have interfered with natural task behavior. However, these reports were used solely for model validation, not training. Thus, our model can reconstruct latent internal states even in tasks without explicit verbalization. Nonetheless, its generalizability should be tested across diverse task settings in future research.
